% !TEX root = ../b-rg.tex
\chapter{Introduction and Context}\label{ch:introduction}

Borlish is a Romance language, the primary language spoken on the island of Borland. It is the only surviving member of the Insular Romance languages.

Throughout this reference we will use `Borlish' to refer to the standard variety of the language, which is based on the educated speech of the capital Damvath.

% --- Voice ---
% The grammar is written in-world, as a description of a real language.
% The only concession to the real world is the use of standard English
% linguistic terminology rather than invented in-universe terms.
%
% --- Planned sections ---
% \section{The language and its speakers}
%   Borland (Istr Boral); Damvath standard; speaker numbers; Borlish abroad
%   (Novomund staddomains?); name of the language (Borallesc).
% \section{Genetic affiliation}
%   Insular Romance (Borlish Latin + British Latin); relation to mainland
%   Romance; Celtic substrate, Old English and Norse superstrata.
% \section{Typological overview}
%   A few pages previewing the whole grammar, with forward references:
%   no gender/number/case on nouns; number via demonstratives (oc/ec);
%   pronouns keep case; rich verb morphology; SVO but object before
%   infinitives; postverbal negation noc; clause-final question particle ou;
%   postnominal adjectives; participial relatives with agentive pronouns.
% \section{Periods of the language}
%   Old Borlish (to c.1250), Middle Borlish, Modern Borlish;
%   the 1870 spelling standard. (Detail goes in the historical appendix.)
% \section{Dialects and registers}
%   Damvath mesolect as the reference variety (cf. Allophony); dialects with
%   syllabic -n̩t -n̩ts (cf. Stress); colloquial vs. literary registers.
% \section{Sources and conventions}
%   Corpus used for examples; notation (\bl, \orth, \phm, \pht, \obs);
%   glossing conventions -> see front-matter Conventions.
